\section{Mountain, Crisis, Witnesses, and Sources}

\subsection{A local event in a continental subsistence crisis}

La Salette-Fallavaux lay above Corps in the southern part of the Diocese of Grenoble, in the Isère Alps. The reported site on Mont Planeau was a high summer pasture, roughly 1,800 meters above sea level. The children were not on pilgrimage or at a devotional gathering. They were tending cattle for farming households near the hamlet of Les Ablandins.

The date, 19 September 1846, fell amid a severe European subsistence crisis. Potato disease, grain failure, rising prices, undernourishment, and insecurity were not images imported into a comfortable society; they belonged to the children's world. The Missionaries' documentary chronology dates the crisis from 1845 and describes the winter of 1846--1847 as famine across France and Europe. The exact local impact varied, and the catastrophe in Ireland must not be treated as interchangeable with conditions in the French Alps. Still, spoiled potatoes, wheat turning to dust, fear of next year's bread, walnuts, grapes, and children dying in famine were immediately intelligible rural realities.

France was still under the July Monarchy. The Revolution of 1848 had not yet occurred, though economic distress, social division, and political instability were building. Religious practice likewise cannot be reduced to a slogan. Catholic institutions and popular devotion were vigorous in many places; anticlericalism, religious indifference, irregular Sunday observance, and uneven catechesis were also real. The public message's complaints about Sunday labor, mockery of Mass, profanity, and Lenten abstinence belong to this concrete pastoral setting.

Historical context performs two necessary tasks and no more. It explains why the reported message used these images and why hearers responded intensely. It does not prove that the children fabricated them; neither does contextual fit prove supernatural origin. A private revelation, if genuine, is received through the language, material fears, and religious grammar of its time.

\subsection{The two children}

Pierre Maximin Giraud was eleven. Born at Corps in 1835, he had lost his mother as an infant, received little schooling or catechesis, and was energetic, distractible, and socially more outgoing than his companion. He was tending cattle for Pierre Selme. Françoise Mélanie Mathieu Calvat was fourteen, approaching fifteen. Born at Corps in 1831 into a large poor household, she had long been hired out for seasonal work, received little formal education, and was tending cattle for Jean-Baptiste Pra. Contemporary descriptions regularly present her as guarded and taciturn.

Both ordinarily spoke the local vernacular and possessed limited French. That matters because the reported speaker began in French, noticed their failure to understand the expression \emph{pommes de terre}, and continued the crop warning in their local speech. Any polished French or English edition is therefore already a translation and an editorial presentation. The children's reported words passed through vernacular retelling, a writer's French, later copies, and modern translation.

They were minors, poor laborers, and newly exposed to intense religious and public scrutiny. Those facts do not make their testimony either automatically pure or automatically worthless. They require child-sensitive historical reading. Repetitive interrogation can test consistency while also shaping vocabulary; admiration can reward elaboration; hostile questioning can produce defensiveness; separation can help test collusion but cannot recreate the event. The record must be evaluated, not romanticized.

\sourceboundary{Maximin and Mélanie are the only reported visual and auditory witnesses to the encounter. Each says the other could see the Lady's lips moving during a private communication but could not hear the words.}{Their agreement and differences, examined across early separate interviews, are the direct human testimony to the event and public message.}{No second observer independently saw or heard the Lady. Witness sincerity does not by itself determine the cause of an experience, and a positive global judgment does not confer verbal inerrancy.}

\subsection{The Pra relation: near the event, not the event itself}

The first written account is dated Sunday, 20 September 1846. That evening Jean-Baptiste Pra, Mélanie's employer, worked with Pierre Selme and Jean Moussier to put the discourse into writing in Mélanie's presence; Maximin had returned to Corps. Its traditional title begins, ``Letter dictated by the Blessed Virgin to two children on the mountain of La Salette-Fallavaux.'' The original has not survived. The text is known through a copy made by Abbé Lagier on 28 February 1847, which he certified as conforming to the original Pra had shown him. Earlier derivative documents agree with important features of it.

That chain gives the relation exceptional value without making it a stenogram. It is next-day testimony, but one child's account is being rendered by adults into French, then known through a later copy. Its awkwardness may preserve proximity; it is not proof that every word came directly from the reported speaker. Jean Stern gives it the first place in his chronological dossier and compares it with the rapidly multiplying early records.

Early documentation soon included the parish priest's and archpriest's reports, interrogations of each child, letters from clergy and lay visitors, reports about the site and claimed favors, and both favorable and hostile press. By October 1846 the existence of a separately spoken communication to each child was present in the dossier. The detailed public narrative became more polished as publications standardized it. The correct method is to identify a stable core and disclose later expansion, not to choose one late seamless text and label it ``the transcript.''

\subsection{The critical dossier}

Jean Stern, M.S., edited the principal modern chronological corpus in three volumes:

\begin{enumerate}
  \item volume 1 covers September 1846 through early March 1847 and contains the first relations and initial reception;
  \item volume 2 continues from late March 1847 through April 1849, including inquiries, pilgrimages, publications, and controversy; and
  \item volume 3 covers 1 May 1849 through 4 November 1854, including the Curé d'Ars episode, the 1851 secret manuscripts and judgment, and Ginoulhiac's act.
\end{enumerate}

The title \emph{Documents authentiques} describes documentary provenance, not a claim that every allegation in every document is true. A hostile letter can be an authentic hostile letter; a witness deposition can authentically record what a person said while leaving the underlying fact disputed. Stern's chronological arrangement is especially important at La Salette because later polemic repeatedly moves a late text backward in time.

Official Shrine and Missionaries pages supply a living custodial narrative. They accurately preserve the event's public center and the continuing devotion. Their brief timelines are pastoral summaries, however, not substitutes for primary acts. In particular, a current Shrine page compresses Ginoulhiac's role into a ``new inquiry'' and attaches the familiar shepherds-and-Church maxim to 1855. The primary sequence is more exact. His substantive formal \emph{Instruction pastorale et Mandement} is dated 4 November 1854. A near-contemporary 1855 printed report says that, on the eve of the anniversary, he described the children's mission as completed and new instruments as taking up the work; his published anniversary allocution then sent the pilgrims themselves out as missionaries of reconciliation. The polished epigram is therefore used here as a later custodial synthesis of an attested pastoral theme, not as a stenographic juridical formula.

\subsection{How the study handles disagreement}

Three distinctions govern the narrative that follows:

\begin{enumerate}
  \item \textbf{Early attestation and later fullness.} A detail found in the next-day relation has a different historical weight from one first printed decades later. Later testimony may still be true, but its date must travel with it.
  \item \textbf{Content and cause.} A historian can establish that a discourse was reported, repeated, and judged credible without directly observing its transcendent cause. The bishop's office made the ecclesial judgment; historical criticism maps the human evidence.
  \item \textbf{Public and personal speech.} The reported commission to communicate the public message is not a license to publish every private communication. A text's sensational interest cannot enlarge the object the bishop judged.
\end{enumerate}

The sources sometimes vary in punctuation, exact wording, the order of minor narrative details, and Maximin's dating of his 1851 manuscript. This study publishes consequential variations and does not manufacture certainty from typographic uniformity. Modern translations are paraphrased except where a short phrase is needed to identify the received message. The underlying French and, where preserved, vernacular witnesses control.
